% Generated draft: requires mathematical and source review.
\documentclass[11pt]{article}
\usepackage[margin=25mm]{geometry}
\usepackage{fontspec}
\setmainfont{Latin Modern Roman}
\usepackage{amsmath,amssymb,mathtools,unicode-math}
\setmathfont{Latin Modern Math}
\usepackage{enumitem,needspace,xurl,hyperref}
\setlength{\parindent}{0pt}
\setlength{\parskip}{5pt}
\setlength{\emergencystretch}{3em}
\newcommand{\N}{\mathbb{N}}
\newcommand{\Z}{\mathbb{Z}}
\newcommand{\Q}{\mathbb{Q}}
\newcommand{\R}{\mathbb{R}}
\newcommand{\C}{\mathbb{C}}
\newcommand{\F}{\mathbb{F}}
\newcommand{\A}{\mathbb{A}}
\newcommand{\PP}{\mathbb{P}}
\newcommand{\E}{\mathbb{E}}
\newcommand{\eps}{\varepsilon}
\newcommand{\dd}{\,\mathrm{d}}
\begin{document}
\section*{UnsolvedMath ID: 30004073}
\textbf{Source problem code:} catalogue entry
\subsection{Definitions and mathematical statement}
Let \(d\) denote the spatial dimension, implicitly a positive integer in the notation \(\R^d\), and let
\[
\frac{2}{d}<p<\frac{4}{d},\qquad \phi\in L^2(\R^d).
\]
Here \(L^2(\R^d)\) denotes the space of square-integrable complex-valued functions, with norm \(\|f\|_2=(\int_{\R^d}|f(x)|^2\,\dd x)^{1/2}\), and \(\Delta=\sum_{\ell=1}^d\partial_{x_\ell}^2\) is the spatial Laplacian. The source denotes a randomization of \(\phi\) by \(\phi^\omega\), but does not specify the randomization or its probability space.

The supplied question asks whether, for every such \(p\) and \(\phi\), for almost every \(\omega\), the solution \(u=u^\omega\) of
\[
\begin{cases}
i\partial_tu-\Delta u=|u|^pu, & x\in\R^d,\\
u(0,x)=\phi^\omega(x)
\end{cases}
\]
admits solitons of the form
\[
u_j(t,x)=e^{it a_j+i x\cdot k_j}\phi_j(x-c_jt)
\]
and a free term, written in the source as \(e^{-it\Delta}\phi\), such that
\[
\lim_{t\to+\infty}
\left\|u(t,\cdot)-\sum_j u_j(t,\cdot)-e^{-it\Delta}\phi\right\|_{L^2(\R^d)}=0.
\]
The operator \(e^{-it\Delta}\) denotes free evolution for \(i\partial_tv-\Delta v=0\).

\textit{Unresolved qualifications.} The displayed assertion preserves the source notation, including its use of the same \(\phi\) for the datum before randomization and for the free profile. It is unclear whether that equality is intended or whether a separately quantified asymptotic profile was intended. The probability law, solution concept, forward global existence assumptions, number of solitons, admissible parameter and profile classes, and any separation conditions are unspecified. If the sum is infinite, its convergence must also be specified. The source calls the \(u_j\) solitons but gives no profile equation or parameter compatibility conditions defining that class. Consequently this is the supplied formulation with its gaps exposed, rather than a fully determined mathematical assertion.
\subsection{Short English statement}
For mass-subcritical NLS with \(2/d<p<4/d\), the question asks whether a solution from randomized \(L^2(\R^d)\) data almost surely approaches a sum of traveling solitons and a free solution in \(L^2\) as \(t\to+\infty\). The randomization, solution and soliton classes, size of the sum, and identity of the free profile require clarification.
\subsection{Sources}
\begin{itemize}
\item Supplied local UnsolvedMath registry record: permanent ID \(30004073\), title \emph{Soliton Resolution for Randomized Mass-Subcritical NLS}, source problem code \texttt{catalogue entry}, in the Oberwolfach Reports Open Problems collection. The supplied proposed year is 2019; no proposer is named.
\item Supplied report citation: \emph{Nonlinear Evolution Equations: Analysis and Numerics} (2020), Oberwolfach Reports, pp.\ 305--405, DOI \texttt{10.4171/owr/2019/5}. The literature-source label is ``Oberwolfach Report 5/2019.'' \url{https://doi.org/10.4171/owr/2019/5}. The background attributes the record to workshop open-problem discussions documented in Oberwolfach Reports. The relationship between the citation year 2020 and the report designation 5/2019 is not explained in the supplied record.
\item Supplied literature source, labeled ``Chatterjee on soliton resolution'': \url{https://arxiv.org/abs/1203.4027}. No fuller bibliographic citation or explanation of its precise relationship to the randomized assertion is supplied.
\item Supplied literature assessment, checked at the exact date \texttt{2026-08-22}: the record labels the problem open and states that the randomized soliton-resolution assertion remains open, while existing work supplies global solutions or statistical and special-case results rather than the stated almost-sure decomposition. This is supplied dated curation derived from the OWR record, not an independently verified literature assessment; its status and cited claims require independent checking.
\item The registry labels the statement \texttt{corrected\_verified} and says its mathematical content was preserved while notation, typography, duplicate clauses, or omitted definitions were repaired. This is supplied editorial metadata, not mathematical verification of the present draft or resolution of the qualifications above.
\end{itemize}
\end{document}
